\begin{textsample}{POS Dim 4 – vem\_ed\_28 – Score 6.00 – vem\_ed\_28/t150.txt}  \label{ex:f4_pos_041}
PCI en Español com conteúdos da Logística Fatec Bebedouro No ensino superior, o intercâmbio virtual ganha um papel de \textbf{destaque} por promover aprendizagem cultural, linguística e tecnológica.
É cada vez mais comum que práticas educacionais mediadas pelas tecnologias contribuam para desenvolver competências específicas e essenciais do Ensino Superior.
No Centro Paula Souza, uma das modalidades de atividade colaborativa que tem ganhado \textbf{espaço} são os Projetos Colaborativos Internacionais (PCIs / Cesu).
Essa metodologia de intercâmbio virtual aproxima professores e alunos de instituições de ensino superior do mundo para discutirem temas da atualidade a partir de conteúdos que envolvem as disciplinas dos cursos tecnológicos.
Em 2024, a Fatec Bebedouro participou da primeira experiência de PCI em parceria com a Uniminuto da Colômbia.
O projeto teve até o momento duas edições, e a temática de interesse comum entre as duas IES foi logística, vinculando as disciplinas de Logística Verde (Fatec Bebedouro) e Gestão da Cadeia de Suprimentos (Uniminuto).
Os grupos mistos de alunos (brasileiros e colombianos), de forma colaborativa, analisaram como a logística sustentável pode contribuir para o desenvolvimento econômico e social para fins da promoção do bem-estar de comunidades locais e globais, estudaram princípios da logística sustentável e sua importância na cadeia de suprimentos moderna e buscaram identificar propostas que contribuam para diminuição do impacto ambiental causado pelos meios de transporte. continuação Essas atividades desenvolvidas por meio da colaboração internacional criam parcerias que potencializam os estudos dos alunos numa perspectiva global, de forma a gerar conhecimento e inovação, que podem ser implementados em seus futuros ambientes de trabalho.
Além disso, a comunicação em idioma estrangeiro, que inicialmente \textbf{parece} ser uma barreira, aos poucos se fortalece.
Especialmente para os alunos fatecanos brasileiros em que o espanhol é componente curricular em Logística, duas horas-aula semanais ao longo de dois semestres.
Ademais, a internacionalização propicia aos professores e estudantes o rompimento de barreiras e a integração da cultura, ensino, pesquisa e a matriz curricular de seus respectivos cursos, de modo que possam criar possibilidades de desenvolvimento dos países e da qualidade de vida das populações. (Stallivieri, 2002).
Propiciar a oportunidade de experiências internacionais abre caminhos para mais competitividade no mercado global.
Esse é um dos papéis “das instituições que buscam o equilíbrio entre as expectativas regionais e nacionais por um lado e os desafios internacionais por outro” (Stallivieri, 2002, local.10).
Quando o intercâmbio virtual acontece entre as disciplinas das áreas técnico-profissionais do curso, o idioma de comunicação pode ser uma barreira.
Por isso, é essencial que o professor de língua estrangeira seja colaborativo para mediar as interações, encorajando e \textbf{apoiando} os alunos na troca de informações relacionadas ao conteúdo em discussão, bem como nas conversas informais que propiciam conhecimento intercultural.
Salomão (2020) considera oportuno que se componham grupos com alunos com maior domínio para atuar como facilitadores sem minimizar o papel dos demais que também contribuem, equilibrando os diferentes papéis na aprendizagem.
Assim, a língua estrangeira vai sendo o motor que possibilita “criar formas de se trabalhar em conjunto na construção de conhecimento ao realizarem uma atividade colaborativa a \textbf{distância}” (Salomão, 2020, p.166), fornecendo andaimes que auxiliam a interação e potencializam as habilidades trazidas pelos alunos.
Outro aspecto importante destacado por Salomão é que o professor de língua estrangeira auxilia o docente da disciplina técnica, nas interações e na construção do projeto.
Estabelecer tal parceria pode ser positivo para a produção de novos materiais, o que torna a aprendizagem mais independente, \textbf{inclusiva} e significativa.
Por se tratar de uma área especializada cuja proposta é discutir dois contextos diferentes (conceitos e aplicações da Logística no contexto brasileiro e colombiano), ainda que houvesse colaboração por parte dos colombianos na comunicação em português, o idioma que se sobressaiu foi a língua espanhola.
Com isso, esse idioma ganhou força na aprendizagem, no caso dos alunos fatecanos brasileiros.
Esse esforço implica um avanço, inclusive no desenvolvimento da língua para fins específicos, pois os alunos passam a ter contato com textos em língua espanhola e a estudar aspectos linguísticos com base em um arcabouço de materiais atrelado a uma especialidade de atuação no mercado de trabalho (Nadin, 2022).
Dessa forma, a aprendizagem se torna mais atrativa e significativa em termos de aplicação na atuação desse aluno em formação.
Referências NADIN, O.
L.
Abordagem Terminológico-Discursiva: pelo resgate da Terminologia e da Terminografia no Ensino e na Aprendizagem de Línguas para Fins Específicos no contexto brasileiro.
Trabalhos em Linguística Aplicada, Campinas, v. 61, p. 97-108, jun. 2022.
STALLIVIERI, L.
O processo de internacionalização nas instituições de ensino superior.
Educação Brasileira - Revista do Conselho de Reitores das Universidades, Brasília, v.24, n.48, p. 35-57, 2002.
Disponível em: https://www.ucs.br/site/midia/arquivos/processo\_internacionalizacao.pdf Acesso em 10 mar. 2025.
SALOMÃO, A.
C.
B.
Intercâmbios virtuais e a internacionalização em casa: reflexões e implicações para a Linguística Aplicada.
São Paulo, Estudos Linguísticos. v.49, n.1, p. 152-174, abr. 2020.
Disponível em: https://revistas.gel.org.br/estudos-linguisticos/article/view/2469/1701.
Acesso em 10 mar. 2025.

% matched lemmas: apoiar, destaque, distância, espaço, inclusivo, parecer
\end{textsample}
